% !TEX root = ../jalase05.tex
%% ============================================================
%%  جلسهٔ ۵ — متغیر تصادفی، تابع جرم احتمال (PMF)،
%%  تابع توزیع تجمعی (CDF)، برنولی، دوجمله‌ای و متغیر نشانگر
%% ============================================================
\chapter{متغیر تصادفی و توزیع‌های گسسته}
\label{ch:s5}

% نمودار میله‌ای PMF:  \stems{رنگ}{k/p, k/p, ...}  (داخل tikzpicture با مقیاس x و y دلخواه)
\providecommand{\stems}[2]{%
  \edef\stemlist{#2}%
  \foreach \k/\p in \stemlist {
    \draw[#1, line width=1.5pt] (\k,0) -- (\k,\p);
    \fill[#1] (\k,\p) circle (1.7pt);
  }}

\sessioncard{شنبه ۱۸ مهر}
{متغیر تصادفی؛ تابع جرم احتمال \lr{(PMF)} و تابع توزیع تجمعی \lr{(CDF)}؛ توزیع‌های برنولی و دوجمله‌ای؛ متغیر نشانگر}
{\lr{Pishro-Nik} بخش‌های \lr{3.1.1} تا \lr{3.1.3}، \lr{3.1.5} و \lr{3.2.1} \,؛\, \lr{Ross} بخش‌های \lr{4.1}، \lr{4.2}، \lr{4.6} و \lr{4.10} \,؛\, \lr{Bertsekas} بخش‌های \lr{2.1}، \lr{2.2} و \lr{3.2}}
{\begin{itemize}
  \item متغیر تصادفی را به‌عنوان «تابعی از فضای نمونه به $\R$» تعریف کنید و پیشامدهایی مثل $\{X=x\}$ و $\{X\in A\}$ را بسازید؛
  \item تابع جرم احتمال یک متغیر گسسته را پیدا و رسم کنید و با آن احتمال هر پیشامدِ مربوط به $X$ را حساب کنید؛
  \item تابع توزیع تجمعی را تعریف کنید، نمودار پله‌ای آن را برای یک متغیر گسسته رسم کنید، خواصش را اثبات کنید و بین \lr{PMF} و \lr{CDF} رفت و برگشت کنید؛
  \item برای توزیع‌های برنولی و دوجمله‌ای «داستان» تولید، \lr{PMF} و شکل نمودار را بدانید و در مسئله تشخیصشان دهید؛
  \item متغیر نشانگرِ یک پیشامد را بسازید، با «جبر نشانگرها» کار کنید و متغیرهای شمارشی را به صورت جمع نشانگرها بنویسید.
\end{itemize}}

%% ------------------------------------------------------------
\section{چرا متغیر تصادفی؟}
%% ------------------------------------------------------------

تا اینجا با سه‌تایی $(\Omega,\F,\P)$ کار کردیم و احتمالِ \emph{پیشامدها} را حساب کردیم. اما در عمل دو مشکل داریم:
\begin{enumerate}
  \item \textbf{فضای نمونه اغلب بسیار پیچیده است.} یک مسابقهٔ فوتبال، وضعیت یک خودرو در حال حرکت، یا ترافیک یک شبکهٔ کامپیوتری را در نظر بگیرید. نوشتن صریح $\Omega$ و احتمال همهٔ زیرمجموعه‌هایش عملاً ناممکن است.
  \item \textbf{ما به همهٔ جزئیات دسترسی نداریم.} فقط چند \emph{عدد} را اندازه می‌گیریم: تعداد گل‌ها، دور موتور، سرعت، تعداد بسته‌های گم‌شده در یک ثانیه، \dots
\end{enumerate}
این اعداد که از روی برآمدِ آزمایش تعیین می‌شوند، \emph{متغیرهای تصادفی}اند. می‌شود آن‌ها را مثل «حسگرهایی» دید که روی یک دنیای تصادفی و پیچیده نصب کرده‌ایم: هر حسگر فقط یک جنبهٔ عددی از برآمد را گزارش می‌دهد.

\begin{figblock}
\begin{tikzpicture}[scale=0.9]
  % فضای نمونه
  \draw[accent, thick, fill=accent!6] plot[smooth cycle, tension=0.8] coordinates
      {(-2.3,0) (-1.6,1.6) (0,2.0) (1.5,1.4) (1.9,-0.3) (0.9,-1.8) (-1.2,-1.7)};
  \node[accent!80!black] at (-1.7,1.9) {$\Omega$};
  \foreach \x/\y in {-1.2/0.6, -0.4/1.1, 0.5/0.7, 1.1/-0.2, 0.2/-0.9, -0.9/-0.7, -0.3/0.1}
     \fill[accent!70!black] (\x,\y) circle (1.3pt);
  \fill[primary] (0.6,0.1) circle (2pt) node[below left, font=\small] {$\omega$};
  % دو محور
  \draw[-{Stealth}, thick, graytext] (5,1.4) -- (10.5,1.4) node[right] {$\R$};
  \draw[-{Stealth}, thick, graytext] (5,-1.2) -- (10.5,-1.2) node[right] {$\R$};
  \draw[setA, thick, smooth, domain=5.3:10, samples=60] plot (\x,{1.4+1.0*exp(-(\x-7.2)^2/0.9)});
  \draw[excolor, thick, smooth, domain=5.3:10, samples=60]
       plot (\x,{-1.2+0.8*exp(-(\x-6.6)^2/0.35)+0.6*exp(-(\x-8.6)^2/0.5)});
  % نگاشت‌ها
  \draw[-{Stealth[length=2.6mm]}, setA, thick] (0.7,0.2) to[bend left=18] node[above, font=\small] {$X(\omega)$} (7.6,1.4);
  \fill[setA] (7.6,1.4) circle (2pt);
  \draw[-{Stealth[length=2.6mm]}, excolor, thick] (0.7,0.0) to[bend right=14] node[below, font=\small] {$Y(\omega)$} (8.4,-1.2);
  \fill[excolor] (8.4,-1.2) circle (2pt);
  \node[font=\small, setA!80!black] at (11.8,2.1) {\rl{توزیع $X$}};
  \node[font=\small, excolor] at (11.8,-0.5) {\rl{توزیع $Y$}};
\end{tikzpicture}

\figcaption[fig:rv-sensors]{متغیرهای تصادفی $X$ و $Y$ دو «اندازه‌گیری» روی یک فضای نمونهٔ پیچیده‌اند. به‌جای کار با کل $\Omega$، با \emph{توزیعِ} این اعداد روی $\R$ کار می‌کنیم.}
\end{figblock}

\begin{intuition}
در یک خودرو، «دور موتور» به ده‌ها عامل بستگی دارد (دندهٔ ماشین، شیب جاده، فشار پدال گاز، \dots). حالت‌های بسیار متفاوتی از خودرو ممکن است دور موتورِ \emph{یکسانی} بدهند؛ پس نگاشت $\omega\mapsto X(\omega)$ معمولاً \key{چند به یک} است و نمی‌توان از روی $X$ به $\omega$ برگشت. یک متغیر تصادفی بخشی از اطلاعات را دور می‌ریزد، اما اطلاعاتی \emph{دربارهٔ} برآمد می‌دهد. با اندازه‌گیری متغیرهای بیشتر (دور موتور، سرعت، دمای موتور) و کشف ارتباط بینشان، تصویر کامل‌تری از سیستم به دست می‌آوریم.
\end{intuition}

%% ------------------------------------------------------------
\section{تعریف متغیر تصادفی}
%% ------------------------------------------------------------

\begin{definition}[label=def:rv]{متغیر تصادفی}
فرض کنید $(\Omega,\F,\P)$ یک فضای احتمال باشد. \term{متغیر تصادفی}{random variable} تابعی است مانند
\[ X:\Omega\to\R \]
که به هر برآمد $\omega\in\Omega$ یک عدد حقیقی $X(\omega)$ نسبت می‌دهد. مجموعهٔ مقادیر ممکن $X$ را \term{برد}{range} آن می‌نامیم و با $R_X=\setc{X(\omega)}{\omega\in\Omega}$ نشان می‌دهیم.
\end{definition}

\begin{note}
\textbf{قرارداد نمادگذاری.} متغیرهای تصادفی را با حروف \emph{بزرگ} ($X,Y,Z,N,\dots$) و مقادیر عددی‌شان را با حروف \emph{کوچک} ($x,y,k,\dots$) نشان می‌دهیم. عبارت «$X=x$» یعنی «متغیر تصادفی $X$ مقدار عددی $x$ را گرفته است».
\end{note}

\begin{definition}[label=def:rv-events]{پیشامدهای ساخته‌شده با متغیر تصادفی}
برای هر عدد $x$ و هر مجموعهٔ $A\subseteq\R$ تعریف می‌کنیم:
\[ \{X=x\}\defeq\setc{\omega\in\Omega}{X(\omega)=x},\qquad \{X\in A\}\defeq\setc{\omega\in\Omega}{X(\omega)\in A}=X^{-1}(A). \]
این‌ها \emph{زیرمجموعه‌هایی از $\Omega$}، یعنی پیشامد، هستند. به جای $\P(\{X=x\})$ می‌نویسیم $\P(X=x)$؛ و ویرگول به معنای «و» است: $\P(X=x,\,Y=y)=\P(\{X=x\}\cap\{Y=y\})$.
\end{definition}

\begin{example}[label=ex:three-coins]{سه پرتاب سکه و یک بازی}
سکهٔ سالمی را سه بار پرتاب می‌کنیم. $X$ را تعداد شیرها بگیرید. در یک بازی، به ازای هر شیر ۱ هزار تومان می‌بریم و به ازای هر خط ۱ هزار تومان می‌بازیم؛ $W$ را سود خالص (به هزار تومان) بگیرید. $R_X$ و $R_W$ را بیابید و $\P(X=2)$ و $\P(W>0)$ را حساب کنید.
\begin{solution}
$\Omega$ هشت برآمد هم‌شانس دارد و دو متغیر تصادفی، دو \emph{تابع} روی آن هستند:
\begin{center}
\renewcommand{\arraystretch}{1.2}
\begin{tabular}{c|cccccccc}
$\omega$ & $HHH$ & $HHT$ & $HTH$ & $THH$ & $HTT$ & $THT$ & $TTH$ & $TTT$\\\hline
$X(\omega)$ & $3$ & $2$ & $2$ & $2$ & $1$ & $1$ & $1$ & $0$\\
$W(\omega)$ & $3$ & $1$ & $1$ & $1$ & $-1$ & $-1$ & $-1$ & $-3$
\end{tabular}
\end{center}
پس $R_X=\{0,1,2,3\}$ و $R_W=\{-3,-1,1,3\}$. پیشامد $\{X=2\}=\{HHT,HTH,THH\}$ است و $\P(X=2)=3/8$. چون $W=X-(3-X)=2X-3$، پیشامد $\{W>0\}$ همان $\{X\ge2\}$ است و $\P(W>0)=\tfrac38+\tfrac18=\tfrac12$.
\end{solution}
\end{example}

\begin{figblock}
\begin{tikzpicture}[yscale=0.62]
  \draw[accent, thick, fill=accent!6, rounded corners=14pt] (-1.1,-4.7) rectangle (1.1,4.7);
  \node[accent!80!black] at (-1.45,4.4) {$\Omega$};
  \foreach \w/\y/\v in {HHH/3.85/3, HHT/2.75/2, HTH/1.65/2, THH/0.55/2,
                        HTT/-0.55/1, THT/-1.65/1, TTH/-2.75/1, TTT/-3.85/0} {
    \node[font=\small] (w\w) at (0,\y) {$\w$};
  }
  \draw[-{Stealth}, thick, graytext] (5.5,-4.6) -- (5.5,4.8) node[above] {$\R$};
  \foreach \v/\y/\pr in {0/-3.3/{1/8}, 1/-1.1/{3/8}, 2/1.1/{3/8}, 3/3.3/{1/8}} {
    \fill[primary] (5.5,\y) circle (2.4pt);
    \node[right, font=\small] at (5.6,\y) {$\v$};
    \node[right, font=\small, accent!80!black] at (6.4,\y) {$\P(X=\v)=\pr$};
  }
  \foreach \w/\c in {HHH/setA} \draw[-{Stealth[length=2mm]}, \c, thick] (w\w.east) -- (5.4,3.3);
  \foreach \w in {HHT,HTH,THH} \draw[-{Stealth[length=2mm]}, excolor, thick] (w\w.east) -- (5.4,1.1);
  \foreach \w in {HTT,THT,TTH} \draw[-{Stealth[length=2mm]}, thmcolor, thick] (w\w.east) -- (5.4,-1.1);
  \draw[-{Stealth[length=2mm]}, warncolor, thick] (wTTT.east) -- (5.4,-3.3);
\end{tikzpicture}

\figcaption[fig:three-coins-map]{متغیر تصادفی «تعداد شیرها» به‌عنوان یک تابع. احتمال هر مقدار، مجموع احتمال برآمدهایی است که به آن مقدار نگاشته می‌شوند.}
\end{figblock}

\begin{warning}
نام «متغیر تصادفی» گمراه‌کننده است: $X$ نه «متغیر» است و نه «تصادفی»؛ یک \emph{تابع کاملاً معین} است. تصادف از انتخاب $\omega$ می‌آید: طبیعت $\omega$ را انتخاب می‌کند و ما فقط $X(\omega)$ را می‌بینیم.
\end{warning}

\subsection{توزیع یک متغیر تصادفی}

هر متغیر تصادفی، قانون احتمال $\P$ روی $\Omega$ را به یک قانون احتمال روی $\R$ «منتقل» می‌کند:
\[ (\Omega,\F,\P)\ \xrightarrow{\ \ X\ \ }\ (\R,\F',\P_X),\qquad \P_X(A)\defeq\P\bigl(X^{-1}(A)\bigr)=\P(X\in A). \]
یعنی زیرمجموعه‌های $\Omega$ به زیرمجموعه‌هایی از $\R$ نگاشته می‌شوند و احتمال‌ها هم همراهشان منتقل می‌شوند.

\begin{proposition}[label=prop:induced]{توزیع القایی}
تابع $\P_X(A)=\P(X\in A)$ یک قانون احتمال روی $\R$ است. این قانون را \term{توزیع}{distribution} متغیر تصادفی $X$ می‌نامیم.
\end{proposition}

\begin{proof}
$\P_X(A)\ge0$ بدیهی است و $\P_X(\R)=\P(X^{-1}(\R))=\P(\Omega)=1$. اگر $A_1,A_2,\dots$ دوبه‌دو مجزا باشند، پیش‌نگاره‌هایشان هم دوبه‌دو مجزا هستند (یک $\omega$ نمی‌تواند به دو مجموعهٔ مجزا نگاشته شود) و $X^{-1}\bigl(\bigcup_iA_i\bigr)=\bigcup_iX^{-1}(A_i)$؛ پس طبق اصل سوم $\P_X\bigl(\bigcup_iA_i\bigr)=\sum_i\P_X(A_i)$.
\end{proof}



%% ------------------------------------------------------------
\section{متغیر تصادفی گسسته و تابع جرم احتمال}
%% ------------------------------------------------------------

\begin{definition}[label=def:discrete-rv]{متغیر تصادفی گسسته}
متغیر تصادفی $X$ \term{گسسته}{discrete} است اگر بردش، $R_X$، متناهی یا نامتناهیِ شمارا باشد. در این صورت می‌توان مقادیرش را فهرست کرد: $R_X=\{x_1,x_2,x_3,\dots\}$.
\end{definition}

\begin{example}{گسسته یا نه؟}
(الف) سکه‌ای را ۱۰۰ بار پرتاب می‌کنیم و $X$ تعداد شیرهاست. (ب) سکه‌ای را آن‌قدر پرتاب می‌کنیم تا اولین شیر بیاید و $Y$ تعداد کل پرتاب‌هاست. (ج) $T$ زمان (به ساعت) از الان تا زلزلهٔ بعدی در تهران است.
\begin{solution}
$R_X=\{0,1,\dots,100\}$ متناهی است و $R_Y=\{1,2,3,\dots\}=\N$ نامتناهیِ شماراست؛ پس هر دو گسسته‌اند. اما $R_T=[0,\infty)$ یک بازه و ناشماراست؛ $T$ گسسته نیست و آن را در جلسه‌های بعد به‌عنوان متغیر \emph{پیوسته} بررسی می‌کنیم.
\end{solution}
\end{example}

برای مدل‌سازی یک متغیر گسسته کافی است بدانیم احتمال هر مقدار چقدر است:
\[ \P(X=x_1)=\,?\qquad \P(X=x_2)=\,?\qquad \dots \]

\begin{definition}[label=def:pmf]{تابع جرم احتمال}
فرض کنید $X$ گسسته باشد. تابع
\[ p_X(x)\defeq\P(X=x),\qquad x\in\R \]
\term{تابع جرم احتمال}{probability mass function, PMF} متغیر $X$ نامیده می‌شود. برای $x\notin R_X$ داریم $p_X(x)=0$.
\end{definition}

به بیان دیگر، \lr{PMF} تابعی است که به تعدادی شمارا از اعداد حقیقی، عددهایی نامنفی نسبت می‌دهد که جمعشان ۱ است:

\begin{theorem}[label=thm:pmf-props]{خواص تابع جرم احتمال}
اگر $X$ متغیری گسسته با برد $R_X$ باشد:
\begin{enumerate}
  \item برای هر $x$: \ $0\le p_X(x)\le1$؛
  \item $\displaystyle\sum_{x\in R_X}p_X(x)=1$؛
  \item برای هر $A\subseteq\R$: \ $\displaystyle\P(X\in A)=\sum_{x\in A\cap R_X}p_X(x)$.
\end{enumerate}
\end{theorem}

\begin{proof}
(۱) از اصل اول و یکنوایی. برای (۲) و (۳): پیشامدهای $\{X=x\}$ برای $x\in R_X$ دوبه‌دو مجزا هستند (هر $\omega$ فقط یک مقدار $X(\omega)$ دارد) و اجتماعشان $\Omega$ است؛ یعنی یک \emph{افراز شمارا} از $\Omega$ می‌سازند. پس طبق اصل سوم، $\sum_{x\in R_X}\P(X=x)=\P(\Omega)=1$. به همین ترتیب $\{X\in A\}$ اجتماع مجزای $\{X=x\}$ها برای $x\in A\cap R_X$ است.
\end{proof}

\begin{note}
\textbf{عکسِ قضیه هم درست است:} اگر $x_1,x_2,\dots$ اعدادی حقیقی و $p_1,p_2,\dots$ اعدادی نامنفی با جمع ۱ باشند، متغیر تصادفی‌ای با $p_X(x_i)=p_i$ وجود دارد (کافی است $\Omega=\{x_1,x_2,\dots\}$، $\P(\{x_i\})=p_i$ و $X(\omega)=\omega$ بگیریم). پس \key{برای مدل‌سازی یک متغیر گسسته، داشتن \lr{PMF} کافی است} و دیگر لازم نیست $\Omega$ را صریحاً بسازیم. آزمایش‌های کاملاً متفاوت (مثلاً «زوج بودن عدد تاس» و «شیر آمدن سکه») ممکن است متغیرهایی با \lr{PMF} یکسان بسازند.
\end{note}

\begin{example}[label=ex:weird-pmf]{مقادیر لازم نیست عدد صحیح باشند}
متغیر $X$ چهار مقدار $-\sqrt2$، $0$، $\pi$ و $10$ را به ترتیب با احتمال‌های $\tfrac13$، $\tfrac14$، $\tfrac16$ و $\tfrac14$ می‌گیرد. آیا این یک \lr{PMF} معتبر است؟ $\P(X>0)$ و $\P(X\in\Z)$ را بیابید.
\begin{solution}
همهٔ احتمال‌ها نامنفی‌اند و $\tfrac13+\tfrac14+\tfrac16+\tfrac14=\tfrac{4+3+2+3}{12}=1$؛ پس معتبر است. طبق خاصیت (۳):
\[ \P(X>0)=p_X(\pi)+p_X(10)=\tfrac16+\tfrac14=\tfrac5{12},\qquad \P(X\in\Z)=p_X(0)+p_X(10)=\tfrac12. \]
\begin{center}
\begin{tikzpicture}[x=0.55cm, y=5.4cm]
  \draw[-{Stealth}, graytext] (-3,0) -- (11.6,0) node[right] {$x$};
  \stems{primary}{-1.414/0.3333, 0/0.25, 3.1416/0.1667, 10/0.25}
  \foreach \x/\lab in {-1.414/-\sqrt2, 0/0, 3.1416/\pi, 10/10} \node[below, font=\small] at (\x,0) {$\lab$};
  \foreach \x/\p/\lab in {-1.414/0.3333/\frac13, 0/0.25/\frac14, 3.1416/0.1667/\frac16, 10/0.25/\frac14}
     \node[above, font=\scriptsize, accent!80!black] at (\x,\p) {$\lab$};
  \node[font=\small] at (-2.4,0.3) {$p_X(x)$};
\end{tikzpicture}
\end{center}
\end{solution}
\end{example}

\begin{example}[label=ex:catan]{مجموع دو تاس و بازی کَتان}
در بازی رومیزی \lr{Catan} هر نوبت دو تاس سالم ریخته می‌شود و زمین‌هایی که شماره‌شان برابر مجموع دو تاس است محصول می‌دهند. بازیکنان باتجربه همیشه کنار زمین‌های ۶ و ۸ خانه می‌سازند و زمین‌های ۲ و ۱۲ را «بی‌ارزش» می‌دانند. چرا؟ \lr{PMF} مجموع $S$ را پیدا کنید.
\begin{solution}
$\Omega$ شامل ۳۶ زوج مرتب هم‌شانس است. تعداد زوج‌هایی با مجموع $s$ برای $s=2,\dots,7$ برابر $s-1$ و برای $s=7,\dots,12$ برابر $13-s$ است؛ به طور فشرده:
\[ p_S(s)=\frac{6-|s-7|}{36},\qquad s=2,3,\dots,12. \]
\begin{center}
\begin{tikzpicture}[x=0.62cm, y=13cm]
  \draw[-{Stealth}, graytext] (1.3,0) -- (12.8,0) node[right] {$s$};
  \stems{primary!45}{2/0.02778, 3/0.05556, 4/0.08333, 5/0.11111, 9/0.11111, 10/0.08333, 11/0.05556, 12/0.02778}
  \stems{accent}{6/0.13889, 8/0.13889}
  \stems{warncolor}{7/0.16667}
  \foreach \s in {2,...,12} \node[below, font=\scriptsize] at (\s,0) {$\s$};
  \foreach \s/\n in {2/1,3/2,4/3,5/4,6/5,7/6,8/5,9/4,10/3,11/2,12/1}
     \node[above, font=\tiny, graytext] at (\s,{\n/36}) {$\frac{\n}{36}$};
\end{tikzpicture}
\end{center}
عدد ۷ محتمل‌ترین مجموع است ($\tfrac16$)؛ اما در کتان، ۷ محصولی نمی‌دهد و «دزد» را فعال می‌کند! بعد از آن، ۶ و ۸ هر کدام با احتمال $\tfrac5{36}$ از بقیه محتمل‌ترند، در حالی که ۲ و ۱۲ فقط با احتمال $\tfrac1{36}$ می‌آیند: زمین ۶ به طور متوسط \emph{پنج برابر} زمین ۱۲ محصول می‌دهد. (روی مهره‌های شمارهٔ بازی هم تعداد نقطه‌ها دقیقاً برابر صورتِ این کسرهاست.)
\end{solution}
\end{example}

\begin{example}[label=ex:heavy-tail]{پیدا کردن ثابت و یک \lr{heavy tail}}
فرض کنید $p_X(k)=\dfrac{c}{k(k+1)}$ برای $k=1,2,3,\dots$. ثابت $c$ را بیابید و $\P(X\ge n)$ را حساب کنید.
\begin{solution}
از تجزیه به کسرهای ساده $\frac1{k(k+1)}=\frac1k-\frac1{k+1}$؛ پس جمع، تلسکوپی است:
\[ \sum_{k=1}^{\infty}\frac{1}{k(k+1)}=\Bigl(1-\tfrac12\Bigr)+\Bigl(\tfrac12-\tfrac13\Bigr)+\cdots=1\quad\Longrightarrow\quad c=1. \]
به همین ترتیب $\P(X\ge n)=\sum_{k\ge n}\bigl(\frac1k-\frac1{k+1}\bigr)=\frac1n$. یعنی با احتمال ۱٪ مقدار $X$ دست‌کم ۱۰۰ و با احتمال ۰٫۱٪ دست‌کم ۱۰۰۰ است! احتمال‌های این توزیع خیلی کُند به صفر می‌روند؛ می‌گوییم چنین توزیعی \emph{\lr{heavy tail}} دارد. (در جلسه‌های بعد خواهیم دید این $X$ حتی «میانگین» متناهی ندارد.)
\end{solution}
\end{example}

%% ------------------------------------------------------------
\section{تابع توزیع تجمعی}
\label{sec:cdf}
%% ------------------------------------------------------------

\lr{PMF} فقط برای متغیرهای گسسته معنا دارد. برای متغیری مثل $T$ (زمان تا زلزلهٔ بعدی) احتمالِ \emph{هر} مقدار مشخص صفر است؛ درست مثل مقاومتِ «دقیقاً ۱ اُهمی» در جلسهٔ ۲. پس $\P(T=t)=0$ برای همهٔ $t$ و \lr{PMF} هیچ اطلاعاتی نمی‌دهد. به ابزاری نیاز داریم که برای \emph{همهٔ} متغیرهای تصادفی کار کند. ایده: به جای احتمالِ تک‌نقطه‌ها، احتمال \key{نیم‌خط‌ها}، یعنی پیشامدهای $\{X\le x\}$، را بدانیم.

\begin{definition}[label=def:cdf]{تابع توزیع تجمعی}
برای هر متغیر تصادفی $X$ (گسسته یا غیر گسسته)، تابع
\[ F_X(x)\defeq\P(X\le x),\qquad x\in\R \]
\term{تابع توزیع تجمعی}{cumulative distribution function, CDF} متغیر $X$ نامیده می‌شود.
\end{definition}

\begin{note}
$F_X$ برای \emph{همهٔ} اعداد حقیقی تعریف شده است، نه فقط برای مقادیر برد. مثلاً اگر $X$ تعداد شیرها در سه پرتاب باشد، $F_X(1.5)=\P(X\le1.5)=\P(X\le1)$؛ چون $X$ بین ۱ و ۱٫۵ مقداری نمی‌گیرد.
\end{note}

برای متغیر گسسته، $F_X(x)$ از «جمع کردن» جرم‌های احتمالِ سمت چپ $x$ به دست می‌آید؛ نام «تجمعی» از همین‌جاست:
\[ F_X(x)=\P(X\le x)=\sum_{\substack{x_i\in R_X\\ x_i\le x}}p_X(x_i). \]
پس تا وقتی $x$ از روی یک مقدارِ برد عبور نکرده، $F_X$ ثابت می‌ماند و وقتی از روی $x_i$ رد می‌شود، به اندازهٔ $p_X(x_i)$ \key{بالا می‌پرد}. نمودار CDF یک متغیر گسسته \emph{پله‌ای} است.

\begin{example}[label=ex:three-coins-cdf]{\lr{CDF} تعداد شیرها}
در مثال \ref{ex:three-coins}، $X$ تعداد شیرها در سه پرتاب سکهٔ سالم است. $F_X$ را بنویسید و رسم کنید.
\begin{solution}
از $p_X(0)=\tfrac18$، $p_X(1)=\tfrac38$، $p_X(2)=\tfrac38$ و $p_X(3)=\tfrac18$ و جمع کردنِ پله به پله:
\[ F_X(x)=\begin{cases}
0 & x<0\\[2pt]
\tfrac18 & 0\le x<1\\[2pt]
\tfrac12 & 1\le x<2\\[2pt]
\tfrac78 & 2\le x<3\\[2pt]
1 & x\ge3
\end{cases} \]
\begin{center}
\begin{tikzpicture}[x=1.45cm, y=3.4cm]
  \draw[-{Stealth}, graytext] (-1.3,0) -- (4.5,0) node[right] {$x$};
  \draw[-{Stealth}, graytext] (-1.2,-0.03) -- (-1.2,1.2) node[above] {$F_X(x)$};
  \foreach \y/\lab in {0.125/\frac18, 0.5/\frac12, 0.875/\frac78, 1/1} {
    \draw[graytext!40, densely dotted] (-1.2,\y) -- (4.3,\y);
    \draw[graytext] (-1.25,\y) -- (-1.15,\y);
    \node[left, font=\scriptsize] at (-1.25,\y) {$\lab$};
  }
  \foreach \x in {0,1,2,3} \node[below, font=\small] at (\x,0) {$\x$};
  % پرش‌ها
  \foreach \x/\a/\b/\lab in {0/0/0.125/\frac18, 1/0.125/0.5/\frac38, 2/0.5/0.875/\frac38, 3/0.875/1/\frac18} {
    \draw[accent, dashed] (\x,\a) -- (\x,\b);
    \node[right, font=\scriptsize, accent!80!black] at (\x,{(\a+\b)/2}) {$+\lab$};
  }
  % پله‌ها
  \draw[primary, very thick] (-1.2,0) -- (0,0);
  \foreach \x/\y in {0/0.125, 1/0.5, 2/0.875} \draw[primary, very thick] (\x,\y) -- ({\x+1},\y);
  \draw[primary, very thick] (3,1) -- (4.3,1);
  \foreach \x/\y in {0/0, 1/0.125, 2/0.5, 3/0.875} \draw[primary, thick, fill=white] (\x,\y) circle (2.2pt);
  \foreach \x/\y in {0/0.125, 1/0.5, 2/0.875, 3/1} \fill[primary] (\x,\y) circle (2.2pt);
\end{tikzpicture}
\end{center}
اندازهٔ هر پرش (نارنجی) دقیقاً مقدار \lr{PMF} در آن نقطه است. در محل پرش، مقدار تابع، مقدارِ \emph{بالایی} است (دایرهٔ توپر)؛ مثلاً $F_X(1)=\tfrac12$ و نه $\tfrac18$، چون $\{X\le1\}$ خودِ نقطهٔ ۱ را هم شامل می‌شود.
\end{solution}
\end{example}

\begin{theorem}[label=thm:cdf-props]{خواص تابع توزیع تجمعی}
برای هر متغیر تصادفی $X$:
\begin{enumerate}
  \item \textbf{نانزولی:} اگر $a<b$ آنگاه $F_X(a)\le F_X(b)$؛
  \item \textbf{حدها در بی‌نهایت:} \ $\displaystyle\lim_{x\to-\infty}F_X(x)=0$ \ و \ $\displaystyle\lim_{x\to+\infty}F_X(x)=1$؛
  \item \textbf{پیوستگی از راست:} برای هر $x$: \ $\displaystyle\lim_{h\downarrow0}F_X(x+h)=F_X(x)$.
\end{enumerate}
\end{theorem}

\begin{proof}
(۱) اگر $a<b$، آنگاه $\{X\le a\}\subseteq\{X\le b\}$ و طبق یکنوایی احتمال $F_X(a)\le F_X(b)$.\\
(۲) چون $F_X$ یکنواست، کافی است حد را روی اعداد صحیح بگیریم. پیشامدهای $\{X\le n\}$ صعودی‌اند و اجتماعشان $\Omega$ است (هر $X(\omega)$ عددی حقیقی و متناهی است)؛ پس طبق پیوستگی احتمال (جلسهٔ ۳) $F_X(n)\to\P(\Omega)=1$. به همین ترتیب $\{X\le-n\}$ نزولی‌اند و اشتراکشان $\varnothing$ است؛ پس $F_X(-n)\to0$.\\
(۳) باز هم کافی است $h=\frac1n$ بگیریم. پیشامدهای $\{X\le x+\frac1n\}$ نزولی‌اند و اشتراکشان $\{X\le x\}$ است؛ پس $F_X(x+\frac1n)\to F_X(x)$.
\end{proof}

\begin{note}
عکس این قضیه هم درست است (اثباتش خارج از این درس است): \key{هر تابعی با این سه خاصیت، \lr{CDF} یک متغیر تصادفی است}. همچنین \lr{CDF} توزیع را \emph{کاملاً} مشخص می‌کند: اگر $F_X=F_Y$ آنگاه $\P(X\in A)=\P(Y\in A)$ برای هر مجموعهٔ $A$.
\end{note}

\begin{proposition}[label=prop:cdf-probs]{محاسبهٔ احتمال با \lr{CDF}}
با نماد $F_X(x^-)\defeq\lim_{h\downarrow0}F_X(x-h)$ (حد از چپ):
\begin{enumerate}
  \item $\P(X>x)=1-F_X(x)$؛
  \item برای $a<b$: \ $\P(a<X\le b)=F_X(b)-F_X(a)$؛
  \item $\P(X<x)=F_X(x^-)$؛
  \item $\P(X=x)=F_X(x)-F_X(x^-)$؛ یعنی \key{احتمال هر نقطه برابر اندازهٔ پرش \lr{CDF} در آن نقطه است}.
\end{enumerate}
\end{proposition}

\begin{proof}
(۱) با متمم‌گیری. (۲) $\{X\le b\}$ اجتماع مجزای $\{X\le a\}$ و $\{a<X\le b\}$ است. (۳) پیشامدهای $\{X\le x-\frac1n\}$ صعودی‌اند و اجتماعشان $\{X<x\}$ است؛ پیوستگی احتمال را به کار ببرید. (۴) $\{X\le x\}$ اجتماع مجزای $\{X<x\}$ و $\{X=x\}$ است.
\end{proof}

\begin{keyformula}
\[ F_X(x)=\sum_{x_i\le x}p_X(x_i)\qquad\Longleftrightarrow\qquad p_X(x)=F_X(x)-F_X(x^-) \]
\rl{برای متغیری با مقادیر صحیح: \ $p_X(k)=F_X(k)-F_X(k-1)$.}
\end{keyformula}

\begin{warning}
برای متغیرهای گسسته، تفاوت «$<$» و «$\le$» \emph{مهم} است. در مثال \ref{ex:three-coins-cdf}:
$\P(X\le2)=F_X(2)=\tfrac78$ ولی $\P(X<2)=F_X(2^-)=\tfrac12$. همچنین
\[ \P(1\le X\le2)=F_X(2)-F_X(1^-)=\tfrac78-\tfrac18=\tfrac34,\qquad\text{\rl{در حالی که}}\qquad F_X(2)-F_X(1)=\tfrac38=\P(1<X\le2). \]
\end{warning}

\begin{example}[label=ex:max-dice]{بزرگ‌ترِ دو تاس: گاهی \lr{CDF} راحت‌تر است}
دو تاس سالم می‌ریزیم و $M$ را بزرگ‌ترِ دو عدد بگیرید. \lr{PMF} متغیر $M$ و $\P(M\ge5)$ را بیابید.
\begin{solution}
شمردن مستقیم $\{M=m\}$ کمی دردسر دارد، ولی $\{M\le m\}$ ساده است: \key{بزرگ‌ترین عدد حداکثر $m$ است اگر و تنها اگر هر دو عدد حداکثر $m$ باشند}. از ۳۶ زوج هم‌شانس، $m\times m$ زوج این ویژگی را دارند:
\[ F_M(m)=\frac{m^2}{36},\qquad m=0,1,\dots,6. \]
حالا \lr{PMF} از تفاضل پله‌ها به دست می‌آید:
\[ p_M(m)=F_M(m)-F_M(m-1)=\frac{m^2-(m-1)^2}{36}=\frac{2m-1}{36},\qquad m=1,\dots,6. \]
(بررسی: $\tfrac{1+3+5+7+9+11}{36}=1$.) و $\P(M\ge5)=1-F_M(4)=1-\tfrac{16}{36}=\tfrac59$.
\end{solution}
\end{example}

\begin{example}[label=ex:cdf-to-pmf]{از \lr{CDF} به \lr{PMF}}
تابع زیر را در نظر بگیرید:
\[ F(x)=\begin{cases}
0 & x<0\\
0.3 & 0\le x<1\\
0.4 & 1\le x<3\\
0.9 & 3\le x<5\\
1 & x\ge5
\end{cases} \]
(الف) آیا $F$ یک \lr{CDF} است؟ (ب) \lr{PMF} متناظر را بیابید. (ج) $\P(1<X\le4)$، $\P(1\le X\le4)$، $\P(X>2)$ و $\P(X=2)$ را حساب کنید.
\begin{solution}
(الف) $F$ نانزولی است، از ۰ شروع و به ۱ ختم می‌شود و در هر نقطهٔ پرش مقدارِ بالایی را دارد (پیوسته از راست)؛ پس طبق قضیهٔ \ref{thm:cdf-props} و نکتهٔ بعدش، یک \lr{CDF} است.\\
(ب) پرش‌ها در $0,1,3,5$ هستند با اندازه‌های $0.3$، $0.1$، $0.5$ و $0.1$؛ پس $R_X=\{0,1,3,5\}$ و $p_X(0)=0.3$، $p_X(1)=0.1$، $p_X(3)=0.5$ و $p_X(5)=0.1$.\\
(ج) $\P(1<X\le4)=F(4)-F(1)=0.9-0.4=0.5$؛ \ $\P(1\le X\le4)=F(4)-F(1^-)=0.9-0.3=0.6$؛ \ $\P(X>2)=1-F(2)=0.6$؛ \ و $\P(X=2)=F(2)-F(2^-)=0$، چون $F$ در ۲ پرشی ندارد.
\end{solution}
\end{example}

\begin{intuition}
\lr{PMF} می‌گوید احتمال «در هر نقطه» چقدر است؛ \lr{CDF} می‌گوید احتمال «تا هر نقطه» چقدر جمع شده است. برای متغیرهای گسسته این دو اطلاعات یکسانی دارند و از هر کدام می‌توان دیگری را ساخت. اما \lr{CDF} کلی‌تر است: در جلسه‌های بعد، برای متغیرهای \emph{پیوسته} نمودار \lr{CDF} به جای پله، یک منحنیِ پیوسته و صعودی از ۰ تا ۱ خواهد بود که هیچ پرشی ندارد؛ هر نقطه احتمال صفر دارد ولی بازه‌ها احتمال مثبت.
\end{intuition}
%% ------------------------------------------------------------
\section{توزیع‌های گسستهٔ مهم}
%% ------------------------------------------------------------

چند الگوی تصادفی آن‌قدر زیاد تکرار می‌شوند که توزیعشان نام دارد. برای هر کدام یک \key{داستان} وجود دارد؛ مهارت اصلی این است که داستان را در صورت مسئله تشخیص دهید، و بعد \lr{PMF} آماده است.

\subsection{توزیع برنولی}

\begin{definition}[label=def:bern]{توزیع برنولی}
متغیر تصادفی $X$ با برد $\{0,1\}$ و
\[ p_X(1)=p,\qquad p_X(0)=1-p\qquad(0\le p\le1) \]
\term{برنولی}{Bernoulli} با پارامتر $p$ نامیده می‌شود و می‌نویسیم $X\sim\mathrm{Bern}(p)$.
\end{definition}

\begin{center}
\begin{tikzpicture}[x=1.6cm, y=2.6cm]
  \draw[-{Stealth}, graytext] (-0.6,0) -- (1.8,0) node[right] {$x$};
  \stems{setA}{0/0.7}
  \stems{accent}{1/0.3}
  \node[below, font=\small] at (0,0) {$0$}; \node[below, font=\small] at (1,0) {$1$};
  \node[left, font=\small] at (0,0.7) {$1-p$}; \node[right, font=\small] at (1,0.3) {$p$};
  \node[font=\small, graytext, anchor=west] at (2.3,0.45) {\rl{داستان: یک آزمایش با دو برآمدِ «موفقیت» (۱) و «شکست» (۰)}};
\end{tikzpicture}
\end{center}

برنولی ساده‌ترین متغیر تصادفی ممکن است، ولی همه‌جا هست: شیر یا خط آمدن یک سکه، رسیدن یا نرسیدن یک بستهٔ شبکه، معیوب بودن یا نبودن یک قطعه، کلیک کردن یا نکردن کاربر روی یک تبلیغ.

\subsection{توزیع دوجمله‌ای}

\textbf{داستان.} $n$ آزمایش برنولی \emph{مستقل} انجام می‌دهیم که در هر کدام احتمال موفقیت $p$ است. $Y$ را \emph{تعداد کل موفقیت‌ها} بگیرید. به زبان متغیرهای تصادفی:
\[ Y=X_1+X_2+\cdots+X_n,\qquad X_1,\dots,X_n\ \overset{\text{i.i.d.}}{\sim}\ \mathrm{Bern}(p). \]
(\lr{i.i.d.} مخفف \lr{independent and identically distributed} یعنی «مستقل و هم‌توزیع» است: هر $X_i$ نتیجهٔ یک آزمایش جداست و همه یک توزیع دارند.)
مقدار $Y$ عددی بین $0$ و $n$ است.

\begin{theorem}[label=thm:binom]{توزیع دوجمله‌ای}
با داستان بالا، $Y\sim\mathrm{Bin}(n,p)$ و
\[ p_Y(k)=\P(Y=k)=\binom nk\,p^k(1-p)^{n-k},\qquad k=0,1,\dots,n. \]
\end{theorem}

\begin{proof}
برآمد آزمایش یک دنبالهٔ $n$تایی از موفقیت ($S$) و شکست ($F$) است. یک دنبالهٔ \emph{مشخص} با دقیقاً $k$ موفقیت، مثلاً $\underbrace{S\cdots S}_{k}\underbrace{F\cdots F}_{n-k}$، به دلیل استقلال آزمایش‌ها احتمال $p^k(1-p)^{n-k}$ دارد؛ و این احتمال به \emph{جای} موفقیت‌ها بستگی ندارد. تعداد چنین دنباله‌هایی برابر تعداد راه‌های انتخاب $k$ جایگاه از $n$ جایگاه، یعنی $\binom nk$ است. این دنباله‌ها پیشامدهای مجزا هستند؛ پس احتمال‌ها جمع می‌شوند.
\end{proof}

\begin{note}
\textbf{چرا «دوجمله‌ای»؟} جمع احتمال‌ها طبق بسط دوجمله‌ای نیوتن برابر ۱ است:
\[ \sum_{k=0}^{n}\binom nk p^k(1-p)^{n-k}=\bigl(p+(1-p)\bigr)^n=1^n=1. \]
احتمال‌های $\mathrm{Bin}(n,p)$ دقیقاً جمله‌های بسط $\bigl(p+(1-p)\bigr)^n$ هستند؛ نام توزیع از همین‌جا آمده است. همچنین $\mathrm{Bin}(1,p)=\mathrm{Bern}(p)$.
\end{note}

\begin{figblock}
\begin{tikzpicture}[x=0.21cm, y=10cm]
  \foreach \sh/\pp/\np/\col/\data in {
     0/0.3/6/setA/{0/0.0008,1/0.0068,2/0.0278,3/0.0716,4/0.1304,5/0.1789,6/0.1916,7/0.1643,8/0.1144,9/0.0654,10/0.0308,11/0.012,12/0.0039,13/0.001,14/0.0002},
     25/0.5/10/excolor/{2/0.0002,3/0.0011,4/0.0046,5/0.0148,6/0.037,7/0.0739,8/0.1201,9/0.1602,10/0.1762,11/0.1602,12/0.1201,13/0.0739,14/0.037,15/0.0148,16/0.0046,17/0.0011,18/0.0002},
     50/0.8/16/accent/{8/0.0001,9/0.0005,10/0.002,11/0.0074,12/0.0222,13/0.0545,14/0.1091,15/0.1746,16/0.2182,17/0.2054,18/0.1369,19/0.0576,20/0.0115}} {
    \begin{scope}[xshift=\sh*0.21cm]
      \draw[graytext] (-0.5,0) -- (20.8,0);
      \foreach \tk in {0,10,20} { \draw[graytext] (\tk,0) -- (\tk,-0.012); \node[below, font=\scriptsize] at (\tk,-0.01) {$\tk$}; }
      \draw[dashed, \col!70!black] (\np,0) -- (\np,0.245);
      \stems{\col}{\data}
      \node[font=\small] at (10,0.29) {$\mathrm{Bin}(20,\,\pp)$};
      \node[font=\scriptsize, \col!70!black, fill=white, inner sep=1pt] at (\np,0.255) {$np=\np$};
    \end{scope}
  }
\end{tikzpicture}

\figcaption[fig:binom-shapes]{\lr{PMF} توزیع دوجمله‌ای با $n=20$. قلهٔ نمودار نزدیک $np$ است و با تغییر $p$ کل نمودار جابه‌جا می‌شود.}
\end{figblock}

\begin{proposition}[label=prop:binom-mode]{شکل نمودار دوجمله‌ای}
اگر $Y\sim\mathrm{Bin}(n,p)$ و $0<p<1$، برای $k=1,\dots,n$:
\[ \frac{p_Y(k)}{p_Y(k-1)}=\frac{(n-k+1)\,p}{k\,(1-p)}. \]
در نتیجه $p_Y(k)$ تا وقتی $k<(n+1)p$ صعودی و پس از آن نزولی است؛ یعنی نمودار «تک‌قله‌ای» است و قله‌اش در $k^*=\lfloor(n+1)p\rfloor$ (نزدیک $np$) قرار دارد.
\end{proposition}

\begin{proof}
با تقسیم دو \lr{PMF}: $\binom nk/\binom n{k-1}=\frac{n-k+1}{k}$ و بقیهٔ عامل‌ها $\frac p{1-p}$ می‌دهند. نسبت بزرگ‌تر از ۱ است اگر و تنها اگر $(n-k+1)p>k(1-p)$، یعنی $k<(n+1)p$.
\end{proof}

\begin{intuition}
وقتی $n$ بزرگ می‌شود: (۱) احتمال دو سرِ نمودار، $\P(Y=0)=(1-p)^n$ و $\P(Y=n)=p^n$، به صفر میل می‌کند؛ (۲) نسبت موفقیت‌ها، $Y/n$، به $p$ نزدیک می‌شود؛ و (۳) تقریباً همهٔ احتمال در بازه‌ای به طول از مرتبهٔ $\sqrt n$ حول $np$ متمرکز است. چون $\sqrt n\ll n$، اگر محور افقی را بر $n$ تقسیم کنیم، نمودار دور $p$ «باریک‌تر و باریک‌تر» می‌شود. این همان پیوند احتمال با فراوانی نسبی (جلسهٔ ۱) است.
\end{intuition}

\begin{example}[label=ex:quiz]{کوییز با جواب‌های شانسی}
کوییزی ۱۰ سؤال چهارگزینه‌ای دارد و برای قبولی دست‌کم ۵ جواب درست لازم است. دانشجویی که هیچ چیز نخوانده، همهٔ سؤال‌ها را مستقل و کاملاً تصادفی جواب می‌دهد. احتمال قبولی‌اش چقدر است؟ محتمل‌ترین تعداد جواب درست چند است؟
\begin{solution}
هر سؤال یک آزمایش برنولی با $p=\tfrac14$ است؛ پس تعداد جواب‌های درست $Y\sim\mathrm{Bin}(10,\tfrac14)$ است:
\[ \P(Y\ge5)=\sum_{k=5}^{10}\binom{10}{k}\Bigl(\frac14\Bigr)^k\Bigl(\frac34\Bigr)^{10-k}\approx0.0584+0.0162+0.0031+0.0004+\cdots\approx0.078. \]
یعنی شانس قبولی با حدس زدن کمتر از ۸٪ است. طبق گزارهٔ \ref{prop:binom-mode}، قله در $\lfloor11/4\rfloor=2$ است: $\P(Y=2)\approx0.282$. احتمال اینکه \emph{هیچ} جواب درستی نداشته باشد هم $(3/4)^{10}\approx0.056$ است.
\end{solution}
\end{example}

\begin{example}[label=ex:repetition]{کد تکرار در کانال نویزی}
روی کانالی که هر بیت را مستقلاً با احتمال $p=0.1$ وارون می‌کند، به جای هر بیت آن را ۳ بار می‌فرستیم ($0\to000$ و $1\to111$) و گیرنده با \emph{رأی اکثریت} تصمیم می‌گیرد. احتمال خطا چقدر است؟ با ۵ بار تکرار چطور؟
\begin{solution}
تعداد بیت‌های خراب‌شده $E\sim\mathrm{Bin}(3,0.1)$ است و رأی اکثریت فقط وقتی اشتباه می‌کند که $E\ge2$:
\[ \P(E\ge2)=\binom32(0.1)^2(0.9)+(0.1)^3=0.027+0.001=0.028. \]
با ۵ تکرار، $E\sim\mathrm{Bin}(5,0.1)$ و خطا یعنی $E\ge3$:
$\P(E\ge3)=10(0.1)^3(0.9)^2+5(0.1)^4(0.9)+(0.1)^5\approx0.0086$.
احتمال خطا از ۱۰٪ به ۲٫۸٪ و سپس به کمتر از ۱٪ رسید؛ بهایش این است که نرخ ارسال اطلاعات ۳ یا ۵ برابر کم شده. یافتن کدهای بهتر از این، موضوع \emph{نظریهٔ کدگذاری} است.
\end{solution}
\end{example}

%% ------------------------------------------------------------
\section{متغیر تصادفی نشانگر}
\label{sec:indicator}
%% ------------------------------------------------------------

هر پیشامد $A$ را می‌توان با یک «چراغ» به عدد تبدیل کرد: اگر $A$ رخ داد چراغ روشن (۱) و وگرنه خاموش (۰) است.

\begin{definition}[label=def:indicator]{متغیر نشانگر}
برای هر پیشامد $A$، \term{متغیر نشانگر}{indicator} $A$ را چنین تعریف می‌کنیم:
\[ \ind_A(\omega)=\begin{cases}1 & \omega\in A\\ 0 & \omega\notin A\end{cases} \]
آنگاه $\ind_A\sim\mathrm{Bern}\bigl(\P(A)\bigr)$، چون $\P(\ind_A=1)=\P(A)$. به عکس، هر متغیر برنولی $X$ نشانگرِ پیشامد $\{X=1\}$ است. گاهی به جای $\ind_A$ می‌نویسیم $\ind\{A\}$.
\end{definition}

\begin{intuition}
متغیر نشانگر پلی بین «پیشامد» و «عدد» است: هر سؤالِ بله/خیر دربارهٔ آزمایش («آیا $A$ رخ داد؟») یک متغیر برنولی می‌سازد. هرچه پیشامد پیچیده‌تر باشد، نشانگرش باز هم فقط ۰ یا ۱ است؛ و همین سادگی آن را به ابزاری قدرتمند برای \key{شمردن} تبدیل می‌کند.
\end{intuition}

\subsection{مثال‌هایی از متغیر نشانگر}

\begin{example}[label=ex:ind-die]{نشانگرها روی یک تاس}
تاس سالمی می‌ریزیم. $A=\{\fa{عدد زوج}\}=\{2,4,6\}$ و $B=\{\fa{عدد دست‌کم ۵}\}=\{5,6\}$. متغیرهای $\ind_A$، $\ind_B$، $\ind_A\ind_B$ و $\ind_A+\ind_B$ را به‌صورت تابعی روی $\Omega$ بنویسید و توزیعشان را بیابید.
\begin{solution}
\begin{center}
\renewcommand{\arraystretch}{1.2}
\begin{tabular}{c|cccccc|l}
$\omega$ & $1$ & $2$ & $3$ & $4$ & $5$ & $6$ & \rl{توزیع}\\\hline
$\ind_A$ & $0$ & $1$ & $0$ & $1$ & $0$ & $1$ & $\mathrm{Bern}(\tfrac12)$\\
$\ind_B$ & $0$ & $0$ & $0$ & $0$ & $1$ & $1$ & $\mathrm{Bern}(\tfrac13)$\\
$\ind_A\ind_B$ & $0$ & $0$ & $0$ & $0$ & $0$ & $1$ & $\mathrm{Bern}(\tfrac16)$\\
$\ind_A+\ind_B$ & $0$ & $1$ & $0$ & $1$ & $1$ & $2$ & \rl{مقادیر $0,1,2$ با احتمال $\tfrac26,\tfrac36,\tfrac16$}
\end{tabular}
\end{center}
حاصل‌ضرب $\ind_A\ind_B$ فقط وقتی ۱ است که \emph{هر دو} رخ دهند؛ پس همان $\ind_{A\cap B}$ است. اما جمع $\ind_A+\ind_B$ دیگر برنولی نیست: این متغیر \key{می‌شمارد که چند تا از پیشامدهای $A$ و $B$ رخ داده‌اند}.
\end{solution}
\end{example}

این مثال نشان می‌دهد که عملیات روی مجموعه‌ها، به حساب روی نشانگرها تبدیل می‌شود:

\begin{proposition}[label=prop:ind-algebra]{جبر نشانگرها}
برای پیشامدهای $A$ و $B$:
\begin{enumerate}
  \item $\ind_{A^c}=1-\ind_A$؛
  \item $\ind_{A\cap B}=\ind_A\,\ind_B$ \ و در نتیجه \ $\ind_A^2=\ind_A$؛
  \item $\ind_{A\cup B}=\ind_A+\ind_B-\ind_A\ind_B$؛ و اگر $A$ و $B$ مجزا باشند، $\ind_{A\cup B}=\ind_A+\ind_B$.
\end{enumerate}
\end{proposition}

\begin{proof}
کافی است تساوی‌ها را برای هر $\omega$ بررسی کنیم. چهار حالت داریم: $\omega$ در هر دو، فقط در $A$، فقط در $B$، یا در هیچ‌کدام. مثلاً برای (۳): در این چهار حالت طرف راست به ترتیب $1+1-1$، $1+0-0$، $0+1-0$ و $0$ است، یعنی دقیقاً $\ind_{A\cup B}$. (۱) و (۲) هم به همین سادگی‌اند و $\ind_A^2=\ind_A\ind_A=\ind_{A\cap A}$.
\end{proof}

\begin{note}
رابطهٔ (۳) همان شمول و عدم شمول (جلسهٔ ۳) است، این بار به زبان اعداد: $\P(A\cup B)=\P(A)+\P(B)-\P(A\cap B)$. در جلسه‌های بعد با «امید ریاضی» خواهیم دید که از چنین تساوی‌هایی بین نشانگرها، تساوی‌های بین احتمال‌ها مستقیماً به دست می‌آید.
\end{note}

\begin{example}[label=ex:ind-parity]{کلید چراغ: نشانگرِ «فرد بودن»}
چراغ یک اتاق در ابتدا خاموش است. $n$ نفر یکی‌یکی از کنار کلید رد می‌شوند و هر کدام، مستقل از بقیه، با احتمال $p$ کلید را می‌زند (یعنی وضعیت چراغ را عوض می‌کند). $A_i$ را پیشامد «نفر $i$ کلید را بزند» و $I_i=\ind_{A_i}$ را نشانگرش بگیرید.
\begin{enumerate}[label=(\alph*)]
  \item[(الف)] \textbf{قبل از هر محاسبه حدس بزنید:} اگر $p=0.01$ و $n=1000$، احتمال اینکه چراغ در پایان روشن باشد بیشتر از $\tfrac12$ است یا کمتر؟ اگر $p=0.9$ باشد چطور؟
  \item[(ب)] نشانگرِ «چراغ در پایان روشن است» را به صورت یک \emph{فرمول جبری} از $I_1,\dots,I_n$ بنویسید؛ بدون حالت‌بندی و بدون عبارت «اگر \dots».
  \item[(ج)] احتمال روشن بودن چراغ را حساب کنید و حدس‌های (الف) را بررسی کنید.
\end{enumerate}
\begin{solution}
\textbf{(ب)} اگر $N=I_1+\cdots+I_n$ تعداد دفعاتِ زدن کلید باشد، چراغ در پایان روشن است اگر و تنها اگر $N$ \emph{فرد} باشد. جمع نشانگرها «تعداد» را می‌دهد، ولی ما فقط \emph{زوج یا فرد بودنِ} آن را می‌خواهیم. ترفند این است که به جای $I_i$ با
\[ S_i\defeq1-2I_i=\begin{cases}+1 & \fa{نفر $i$ کلید را نزند}\\ -1 & \fa{نفر $i$ کلید را بزند}\end{cases} \]
کار کنیم. اگر «خاموش» را با $+1$ و «روشن» را با $-1$ نشان دهیم، زدن کلید یعنی \key{ضرب در $-1$}؛ پس وضعیت نهایی چراغ
\[ S_1S_2\cdots S_n=(-1)^{N} \]
است: اگر $N$ زوج باشد $+1$ و اگر فرد باشد $-1$. حالا کافی است $\{+1,-1\}$ را به $\{0,1\}$ برگردانیم:
\[ \boxed{\ \ind\{\fa{چراغ روشن}\}=\ind\{N\ \fa{فرد}\}=\frac{1-\prod_{i=1}^{n}(1-2I_i)}{2}\ } \]
\medskip

\textbf{(ج)} هر نفر یک آزمایش برنولی مستقل با احتمال موفقیت $p$ است؛ پس طبق داستان دوجمله‌ای $N\sim\mathrm{Bin}(n,p)$. همان علامتِ $(-1)^N$ از قسمت (ب) را در \lr{PMF} ضرب کنید و بسط دوجمله‌ای نیوتن را به کار ببرید:
\[ \P(N\ \fa{زوج})-\P(N\ \fa{فرد})=\sum_{k=0}^{n}(-1)^k\binom nk p^k(1-p)^{n-k}=\bigl((1-p)-p\bigr)^n=(1-2p)^n. \]
از طرف دیگر $\P(N\ \fa{زوج})+\P(N\ \fa{فرد})=1$. با کم کردن دو معادله:
\[ \P(\fa{چراغ روشن})=\frac{1-(1-2p)^n}{2}. \]
(بررسی: برای $n=1$ جواب $p$ است و برای $n=2$ جواب $2p(1-p)$، یعنی «دقیقاً یک نفر کلید بزند».)

\textbf{بررسی حدس‌ها.}
\begin{itemize}
  \item $p=0.01$، $n=1000$: \ $(0.98)^{1000}\approx e^{-20.2}\approx2\times10^{-9}$؛ پس احتمال \emph{تقریباً دقیقاً} $\tfrac12$ است، با اینکه به طور متوسط فقط حدود ۱۰ نفر کلید را می‌زنند.
  \item $p=0.9$: \ $1-2p=-0.8$ منفی است و جواب حول $\tfrac12$ \key{نوسان} می‌کند: برای $n=2$ برابر $\tfrac{1-0.64}{2}=0.18$ و برای $n=3$ برابر $\tfrac{1+0.512}{2}=0.756$. وقتی تقریباً همه کلید را می‌زنند، زوج یا فرد بودن $n$ تعیین‌کننده است.
  \item $p=\tfrac12$: جواب برای \emph{هر} $n\ge1$ دقیقاً $\tfrac12$ است. و $p=1$: چراغ به طور قطعی برای $n$ فرد روشن و برای $n$ زوج خاموش است.
\end{itemize}
در حالت کلی برای هر $0<p<1$ داریم $|1-2p|<1$؛ پس با بزرگ شدن $n$ احتمال به $\tfrac12$ میل می‌کند: چراغ بعد از تعداد زیادی تغییرِ تصادفی، وضعیت اولیه‌اش را «فراموش» می‌کند.
\end{solution}
\end{example}

%% ------------------------------------------------------------
\section*{جمع‌بندی جلسهٔ پنجم}
\addcontentsline{toc}{section}{خلاصه و تمرین‌ها}
%% ------------------------------------------------------------

\begin{summary}
\begin{itemize}
  \item \textbf{متغیر تصادفی} تابعی $X:\Omega\to\R$ است؛ «اندازه‌گیری» عددی از یک آزمایش تصادفی. پیشامدها: $\{X\in A\}=X^{-1}(A)$. توزیع $X$ یعنی $\P_X(A)=\P(X\in A)$ که خودش یک قانون احتمال روی $\R$ است.
  \item $X$ \textbf{گسسته} است اگر بردش شمارا باشد. \textbf{\lr{PMF}:} $p_X(x)=\P(X=x)$؛ نامنفی با جمع ۱، و $\P(X\in A)=\sum_{x\in A}p_X(x)$. داشتن \lr{PMF} برای مدل‌سازی کافی است.
  \item \textbf{\lr{CDF}:} $F_X(x)=\P(X\le x)$ برای هر متغیر تصادفی تعریف می‌شود؛ نانزولی، از ۰ تا ۱، و پیوسته از راست است. $\P(a<X\le b)=F_X(b)-F_X(a)$ و احتمال هر نقطه برابر اندازهٔ پرش $F_X$ در آن است. برای متغیر گسسته نمودارش پله‌ای است و ارتفاع پله‌ها همان \lr{PMF} است.
  \item \textbf{برنولی:} یک آزمایش موفقیت/شکست؛ $p_X(1)=p$.
  \item \textbf{دوجمله‌ای:} تعداد موفقیت‌ها در $n$ آزمایش مستقل؛ $\binom nkp^k(1-p)^{n-k}$. نمودار تک‌قله‌ای با قله نزدیک $np$ و پهنای از مرتبهٔ $\sqrt n$.
  \item \textbf{نشانگر:} $\ind_A\sim\mathrm{Bern}(\P(A))$؛ عملیات مجموعه‌ای به حساب تبدیل می‌شود: $\ind_{A^c}=1-\ind_A$ و $\ind_{A\cap B}=\ind_A\ind_B$. متغیرهای شمارشی را می‌توان جمع نشانگرها نوشت.
\end{itemize}
\end{summary}

\subsection*{تمرین‌های پیشنهادی}

\begin{exercise}
سکهٔ سالمی را آن‌قدر پرتاب می‌کنیم تا اولین شیر بیاید، اما حداکثر ۴ بار. $N$ را تعداد پرتاب‌ها بگیرید. $R_N$ و \lr{PMF} آن را بیابید و بررسی کنید جمع احتمال‌ها ۱ است. سپس $F_N$ را بنویسید و رسم کنید.
\answer{$R_N=\{1,2,3,4\}$ با احتمال‌های $\tfrac12,\tfrac14,\tfrac18,\tfrac18$؛ \ $F_N$ در $1,2,3,4$ به مقادیر $\tfrac12,\tfrac34,\tfrac78,1$ می‌پرد.}
\end{exercise}

\begin{exercise}
اگر $p_X(k)=c\,(1/3)^k$ برای $k=0,1,2,\dots$، ثابت $c$ را بیابید. احتمال اینکه $X$ زوج باشد چقدر است؟ $F_X(n)$ را برای عدد صحیح $n\ge0$ حساب کنید.
\answer{$c=\tfrac23$، \ $\P(\fa{$X$ زوج})=\tfrac34$ \ و \ $F_X(n)=1-(1/3)^{n+1}$}
\end{exercise}

\begin{exercise}[\lr{Ross}]
از ظرفی با ۲۰ مهرهٔ شماره‌دار (۱ تا ۲۰)، ۳ مهره بدون جایگذاری برمی‌داریم. $X$ بزرگ‌ترین شمارهٔ برداشته‌شده است. $p_X(i)$ و $\P(X\ge17)$ را بیابید. (راهنمایی: مانند مثال \ref{ex:max-dice}، ابتدا $F_X(i)=\P(X\le i)$ را حساب کنید.)
\answer{$F_X(i)=\binom{i}{3}\big/\binom{20}{3}$، \ $p_X(i)=\binom{i-1}{2}\big/\binom{20}{3}$ برای $i=3,\dots,20$ \ و \ $\P(X\ge17)=1-\binom{16}3\big/\binom{20}3\approx0.509$}
\end{exercise}

\begin{exercise}
دو تاس سالم می‌ریزیم و $L$ را کوچک‌ترِ دو عدد بگیرید. ابتدا $\P(L>m)$ را حساب کنید و از آن $F_L$ و $p_L$ را به دست آورید.
\answer{$\P(L>m)=\bigl(\tfrac{6-m}{6}\bigr)^2$، پس $F_L(m)=1-\tfrac{(6-m)^2}{36}$ و $p_L(m)=\tfrac{13-2m}{36}$ برای $m=1,\dots,6$.}
\end{exercise}

\begin{exercise}
کدام‌یک از توابع زیر \lr{CDF} است؟ برای هر کدام که نیست، خاصیتِ نقض‌شده را بگویید.
(الف) $F(x)=0$ برای $x<0$ و $F(x)=1-\tfrac1{2^{\lfloor x\rfloor+1}}$ برای $x\ge0$. \
(ب) $F(x)=0$ برای $x\le0$ و $F(x)=1$ برای $x>0$. \
(ج) $F(x)=\tfrac12$ برای $x<1$ و $F(x)=1$ برای $x\ge1$. \
(د) $F(x)=0$ برای $x<0$، $F(x)=0.6$ برای $0\le x<1$، $F(x)=0.4$ برای $1\le x<2$ و $F(x)=1$ برای $x\ge2$.
\answer{(الف) بله؛ پرش‌ها در $k=0,1,2,\dots$ با اندازهٔ $p_X(k)=2^{-(k+1)}$؛ (ب) نه، در ۰ پیوسته از راست نیست؛ (ج) نه، در $-\infty$ به ۰ میل نمی‌کند؛ (د) نه، نزولی می‌شود.}
\end{exercise}

\begin{exercise}[مسئلهٔ شوالیه دومِره]
کدام محتمل‌تر است: دست‌کم یک ۶ در ۴ بار ریختن یک تاس، یا دست‌کم یک جفت‌شش در ۲۴ بار ریختن دو تاس؟ (همان پرسشی که در قرن هفدهم پاسکال و فرما را به نامه‌نگاری واداشت؛ جلسهٔ ۱.) هر دو را به زبان یک متغیر دوجمله‌ای بیان کنید.
\answer{$1-(5/6)^4\approx0.518$ \ در برابر \ $1-(35/36)^{24}\approx0.491$؛ اولی محتمل‌تر است.}
\end{exercise}

\begin{exercise}[پیش‌فروش بیش از ظرفیت]
یک هواپیما ۵۰ صندلی دارد ولی شرکت ۵۲ بلیت می‌فروشد، چون هر مسافر مستقلاً با احتمال $0.9$ در فرودگاه حاضر می‌شود. احتمال اینکه مسافری بدون صندلی بماند چقدر است؟
\answer{$\P(X\ge51)=52(0.9)^{51}(0.1)+(0.9)^{52}\approx0.028$}
\end{exercise}

\begin{exercise}
اگر $Y\sim\mathrm{Bin}(n,p)$، نشان دهید $n-Y\sim\mathrm{Bin}(n,1-p)$. یک بار با \lr{PMF} و یک بار با نوشتن $Y$ به صورت جمع نشانگرها استدلال کنید.
\end{exercise}

\begin{exercise}
با جبر نشانگرها نشان دهید
\[ \ind_{A\cup B\cup C}=1-(1-\ind_A)(1-\ind_B)(1-\ind_C), \]
حاصل‌ضرب را باز کنید و آن را با فرمول شمول و عدم شمول برای سه پیشامد مقایسه کنید. همچنین نشان دهید نشانگر تفاضل متقارن $A\triangle B$ برابر $(\ind_A-\ind_B)^2$ است.
\end{exercise}

\begin{exercise}
تاسی سالم را ۱۰ بار می‌ریزیم. $Y$ را تعداد دفعاتی بگیرید که عدد ۱ یا ۲ آمده است. $Y$ را به صورت جمع نشانگرها بنویسید، توزیعش را مشخص کنید و $\P(Y=3)$ را حساب کنید.
\answer{$Y=\sum_{t=1}^{10}\ind\{\fa{پرتاب $t$ برابر ۱ یا ۲}\}\sim\mathrm{Bin}(10,\tfrac13)$ \ و \ $\P(Y=3)=\binom{10}3\tfrac{2^7}{3^{10}}\approx0.260$}
\end{exercise}

\begin{exercise}
در یک کلاس، ۴ دانشجو هر کدام یک تاس سالم می‌ریزند. $Z$ را تعداد \emph{زوج‌هایی} از دانشجوها بگیرید که عدد یکسان آورده‌اند. $Z$ را به صورت جمع ۶ نشانگر بنویسید. هر نشانگر چه توزیعی دارد؟ آیا $Z\sim\mathrm{Bin}(6,\tfrac16)$؟ (راهنمایی: $\P(Z=5)$ را بررسی کنید.)
\answer{$Z=\sum_{i<j}\ind\{\fa{دانشجوی $i$ و $j$ عدد یکسان}\}$، هر کدام $\mathrm{Bern}(\tfrac16)$؛ اما $Z$ دوجمله‌ای نیست: اگر ۵ زوج یکسان باشند زوج ششم هم یکسان است، پس $\P(Z=5)=0$.}
\end{exercise}

\begin{exercise}\hard
سکهٔ سالمی را $n$ بار پرتاب می‌کنیم. «تغییر» یعنی جایی که پرتاب $i$ با پرتاب $i+1$ متفاوت است. $Z$ را تعداد تغییرها بگیرید (مثلاً در $HHTHT$ سه تغییر داریم). $Z$ را جمع $n-1$ نشانگر بنویسید و نشان دهید $Z\sim\mathrm{Bin}(n-1,\tfrac12)$. (راهنمایی: نشان دهید هر الگوی تغییر/عدم‌تغییر دقیقاً با ۲ دنبالهٔ پرتاب سازگار است.) اگر سکه سالم نباشد چه؟
\answer{هر یک از $2^{n-1}$ الگوی تغییر احتمال $2/2^n=2^{-(n-1)}$ دارد؛ پس نشانگرها مثل $n-1$ آزمایش مستقل با $p=\tfrac12$ رفتار می‌کنند. برای سکهٔ ناسالم، $Z$ در حالت کلی دوجمله‌ای نیست.}
\end{exercise}
